% section: 複素数の極形式と指数表示

\section{複素数の極形式と指数表示}

複素数を極形式で表すと,
\[
  z
  =
  \rho(
    \cos\theta+i\sin\theta
  )
\]
と書くことができる.

このとき,
\[
  \cos\theta+i\sin\theta
\]
のような複素数を何乗かしたときに,
偏角がどのように変化するかを考える.

まず,三角関数の加法定理を用いる.


\subsection{極形式の積}

  任意の実数$\alpha,\beta$について,
  \begin{align*}
    &(\cos\alpha+i\sin\alpha)
    (\cos\beta+i\sin\beta)\\
    &=
    \cos\alpha\cos\beta
    -\sin\alpha\sin\beta\\
    &\qquad
    +i(
      \sin\alpha\cos\beta
      +\cos\alpha\sin\beta
    ).
  \end{align*}

  加法定理より,
  \[
    \cos\alpha\cos\beta
    -\sin\alpha\sin\beta
    =
    \cos(\alpha+\beta)
  \]
  および
  \[
    \sin\alpha\cos\beta
    +\cos\alpha\sin\beta
    =
    \sin(\alpha+\beta)
  \]
  であるから,
  \[
    (\cos\alpha+i\sin\alpha)
    (\cos\beta+i\sin\beta)
    =
    \cos(\alpha+\beta)
    +i\sin(\alpha+\beta)
  \]
  を得る.

  つまり,
  \[
    \cos\alpha+i\sin\alpha
  \]
  と
  \[
    \cos\beta+i\sin\beta
  \]
  を掛けると,絶対値は1のままで偏角が
  \[
    \alpha+\beta
  \]
  となる.

  この性質を繰り返し用いることで,
  複素数の累乗について重要な公式が得られる.


\subsection{ド・モアブルの定理}

  任意の整数$n$について,
  \[
    (\cos\theta+i\sin\theta)^n
    =
    \cos(n\theta)+i\sin(n\theta)
  \]
  が成り立つ.

  まず,$n$が正の整数の場合を考える.

  $n=1$のとき,
  \[
    (\cos\theta+i\sin\theta)
    =
    \cos\theta+i\sin\theta
  \]
  であるから,成り立つ.

  いま,
  \[
    (\cos\theta+i\sin\theta)^n
    =
    \cos(n\theta)+i\sin(n\theta)
  \]
  が成り立つと仮定する.

  このとき,
  \begin{align*}
    &(\cos\theta+i\sin\theta)^{n+1}\\
    &=
    \{
      \cos(n\theta)+i\sin(n\theta)
    \}
    (\cos\theta+i\sin\theta)\\
    &=
    \cos((n+1)\theta)
    +i\sin((n+1)\theta)
  \end{align*}
  となる.

  よって,数学的帰納法により,
  \[
    (\cos\theta+i\sin\theta)^n
    =
    \cos(n\theta)+i\sin(n\theta)
  \]
  がすべての正の整数$n$について成り立つ.

  $n=0$のときは,
  \[
    (\cos\theta+i\sin\theta)^0
    =
    1
    =
    \cos0+i\sin0
  \]
  である.

  また,
  \[
    \cos(-\theta)=\cos\theta,
    \qquad
    \sin(-\theta)=-\sin\theta
  \]
  より,
  \begin{align*}
    &(\cos\theta+i\sin\theta)
    (\cos(-\theta)+i\sin(-\theta))\\
    &=
    \cos0+i\sin0\\
    &=
    1.
  \end{align*}

  したがって,
  \[
    (\cos\theta+i\sin\theta)^{-1}
    =
    \cos(-\theta)+i\sin(-\theta)
  \]
  である.

  これを用いれば,負の整数についても同じ公式が成り立つ.

  したがって,
  \[
    (\cos\theta+i\sin\theta)^n
    =
    \cos(n\theta)+i\sin(n\theta)
    \qquad(n\in\mathbb Z)
  \]
  が成り立つ.

  これをド・モアブルの定理という.

  例えば,
  \[
    \left(
      \cos\frac{\pi}{6}
      +i\sin\frac{\pi}{6}
    \right)^6
    =
    \cos\pi+i\sin\pi
    =
    -1
  \]
  となる.

  このように,複素数を極形式で表すと,累乗を考えるときには絶対値の累乗と偏角の$n$倍を考えればよいことが分かる.


\subsection{複素数の指数表示}

  次に,実数の指数関数で用いられる
  \[
    e^x
    =
    \lim_{n\to\infty}
    \left(
      1+\frac{x}{n}
    \right)^n
  \]
  という極限表示を複素数に拡張する.

  ここでは,
  \[
    e^{i\theta}
    :=
    \lim_{n\to\infty}
    \left(
      1+\frac{i\theta}{n}
    \right)^n
  \]
  と定める.

  この極限を評価するため,
  \[
    1+\frac{i\theta}{n}
  \]
  を極形式に直す.

  その絶対値は
  \[
    r_n
    =
    \sqrt{
      1+\frac{\theta^2}{n^2}
    }
  \]
  である.

  また,$\phi_n$を
  \[
    \tan\phi_n
    =
    \frac{\theta}{n}
  \]
  を満たす角とする.

  このとき,
  \[
    \cos\phi_n
    =
    \frac{1}{r_n}
  \]
  および
  \[
    \sin\phi_n
    =
    \frac{\theta}{nr_n}
  \]
  であるから,
  \[
    1+\frac{i\theta}{n}
    =
    r_n(
      \cos\phi_n+i\sin\phi_n
    )
  \]
  と表すことができる.

  よって,
  \begin{align*}
    \left(
      1+\frac{i\theta}{n}
    \right)^n
    &=
    r_n^n
    (\cos\phi_n+i\sin\phi_n)^n\\
    &=
    r_n^n
    \{
      \cos(n\phi_n)
      +i\sin(n\phi_n)
    \}.
  \end{align*}

  最後の等式にはド・モアブルの定理を用いた.

  したがって,この極限を求めるためには,
  \[
    r_n^n
  \]
  と
  \[
    n\phi_n
  \]
  の極限を調べればよい.


\subsection{極限の計算}

  まず,
  \[
    r_n^n
    =
    \left(
      1+\frac{\theta^2}{n^2}
    \right)^{n/2}
  \]
  について考える.

  自然対数の基本極限
  \[
    \lim_{x\to0}
    \frac{\log(1+x)}{x}
    =
    1
  \]
  を用いると,
  \begin{align*}
    \log(r_n^n)
    &=
    \frac n2
    \log
    \left(
      1+\frac{\theta^2}{n^2}
    \right)\\
    &=
    \frac{\theta^2}{2n}
    \frac{
      \log
      \left(
        1+\frac{\theta^2}{n^2}
      \right)
    }{
      \frac{\theta^2}{n^2}
    }.
  \end{align*}

  $n\to\infty$のとき,
  \[
    \frac{\theta^2}{2n}\to0
  \]
  であり,
  \[
    \frac{
      \log
      \left(
        1+\frac{\theta^2}{n^2}
      \right)
    }{
      \frac{\theta^2}{n^2}
    }
    \to1
  \]
  であるから,
  \[
    \log(r_n^n)\to0.
  \]

  よって,
  \[
    r_n^n\to1
  \]
  である.

  次に,
  \[
    n\phi_n
  \]
  の極限を考える.

  $\tan\phi_n=\dfrac{\theta}{n}$より,
  \[
    \tan\phi_n\to0
  \]
  であるから,
  \[
    \phi_n\to0
  \]
  である.

  また,
  \[
    \tan x
    =
    \frac{\sin x}{\cos x}
  \]
  および,
  \[
    \lim_{x\to0}\frac{\sin x}{x}=1,
    \qquad
    \lim_{x\to0}\cos x=1
  \]
  より,
  \[
    \lim_{x\to0}\frac{\tan x}{x}=1
  \]
  である.

  したがって,
  \[
    \lim_{x\to0}\frac{x}{\tan x}=1
  \]
  であるから,
  \begin{align*}
    n\phi_n
    &=
    \theta
    \frac{\phi_n}{\tan\phi_n}\\
    &\to
    \theta.
  \end{align*}

  よって,
  \[
    n\phi_n\to\theta
  \]
  が成り立つ.


\subsection{オイラーの公式}

  以上より,
  \[
    r_n^n\to1
  \]
  および
  \[
    n\phi_n\to\theta
  \]
  である.

  $\sin x,\cos x$は連続関数であるから,
  \[
    \cos(n\phi_n)\to\cos\theta
  \]
  および
  \[
    \sin(n\phi_n)\to\sin\theta
  \]
  となる.

  したがって,
  \begin{align*}
    e^{i\theta}
    &=
    \lim_{n\to\infty}
    \left(
      1+\frac{i\theta}{n}
    \right)^n\\
    &=
    \lim_{n\to\infty}
    r_n^n
    \{
      \cos(n\phi_n)
      +i\sin(n\phi_n)
    \}\\
    &=
    \cos\theta+i\sin\theta.
  \end{align*}

  よって,
  \[
    e^{i\theta}
    =
    \cos\theta+i\sin\theta
  \]
  が成り立つ.

  これをオイラーの公式という.

  したがって,
  \[
    \cos\theta+i\sin\theta
    =
    e^{i\theta}
  \]
  と書くことができる.

  これにより,複素数
  \[
    z
    =
    \rho(
      \cos\theta+i\sin\theta
    )
  \]
  は,
  \[
    z
    =
    \rho e^{i\theta}
  \]
  と表せる.


\subsection{ド・モアブルの定理の指数表示}

  オイラーの公式を用いると,
  ド・モアブルの定理は指数関数を使ってさらに簡潔に表すことができる.

  実際,
  \[
    \cos\theta+i\sin\theta
    =
    e^{i\theta}
  \]
  より,
  \begin{align*}
    (\cos\theta+i\sin\theta)^n
    &=
    (e^{i\theta})^n\\
    &=
    e^{in\theta}\\
    &=
    \cos(n\theta)+i\sin(n\theta).
  \end{align*}

  また,
  \[
    z
    =
    \rho e^{i\theta}
  \]
  と表せば,
  \begin{align*}
    z^n
    &=
    \rho^n e^{in\theta}\\
    &=
    \rho^n
    \{
      \cos(n\theta)+i\sin(n\theta)
    \}.
  \end{align*}

  したがって,
  \[
    z^n
    =
    \rho^n
    \{
      \cos(n\theta)+i\sin(n\theta)
    \}
    \qquad(n\in\mathbb Z)
  \]
  となる.

  ここでは,極形式における「絶対値は$n$乗し,偏角は$n$倍する」という性質が,
  指数表示によって簡潔に表されている.


\subsection{実数指数への拡張}

  これまでのド・モアブルの定理では,指数$n$は整数であった.

  ここで,指数を一般の実数にまで広げることを考える.

  まず,複素数
  \[
    z
    =
    \rho(
      \cos\theta+i\sin\theta
    )
    =
    \rho e^{i\theta}
    \qquad(\rho>0)
  \]
  に対して,実数$x$を指数とする場合を考える.

  実数の指数関数では,
  \[
    \rho^x
  \]
  が定義されているので,
  \[
    z^x
  \]
  を
  \[
    z^x
    :=
    \rho^x e^{ix\theta}
  \]
  と定める.

  オイラーの公式より,
  \[
    e^{ix\theta}
    =
    \cos(x\theta)+i\sin(x\theta)
  \]
  であるから,
  \[
    z^x
    =
    \rho^x
    \{
      \cos(x\theta)+i\sin(x\theta)
    \}
  \]
  となる.

  したがって,
  \[
    z^x
    =
    \rho^x
    \{
      \cos(x\theta)+i\sin(x\theta)
    \}
    \qquad(x\in\mathbb R)
  \]
  と表すことができる.

  $x$が整数の場合には,
  これはすでに示したド・モアブルの定理と一致する.

  また,$x=\dfrac{m}{n}$が有理数の場合には,
  \[
    z^{m/n}
    =
    \rho^{m/n}
    \left\{
      \cos\frac{m\theta}{n}
      +i\sin\frac{m\theta}{n}
    \right\}
  \]
  である.

  この数を$n$乗すると,
  \begin{align*}
    \left(
      z^{m/n}
    \right)^n
    &=
    \rho^m
    \left(
      \cos\frac{m\theta}{n}
      +i\sin\frac{m\theta}{n}
    \right)^n\\
    &=
    \rho^m
    \{
      \cos(m\theta)
      +i\sin(m\theta)
    \}\\
    &=
    z^m.
  \end{align*}

  したがって,$z^{m/n}$は$z^m$の$n$乗根となっており,
  この定義は有理数の場合にも通常の指数の意味と整合している.

  以上から,ド・モアブルの定理は,
  \[
    z^x
    =
    \rho^x
    \{
      \cos(x\theta)+i\sin(x\theta)
    \}
    \qquad(x\in\mathbb R)
  \]
  という形に拡張することができる.

  ただし,複素数の偏角は一般には一意ではなく,
  \[
    \theta,\quad
    \theta+2\pi,\quad
    \theta+4\pi,\quad\ldots
  \]
  はすべて同じ複素数を表す.

  このため,一般の複素数の実数乗を考えるときには,
  どの偏角を用いるかをあらかじめ固定しておく必要がある.

  本書では,
  \[
    z
    =
    \rho(
      \cos\theta+i\sin\theta
    )
  \]
  において$\theta$を一つ固定し,
  \[
    z^x
    =
    \rho^x
    \{
      \cos(x\theta)+i\sin(x\theta)
    \}
  \]
  と考えることにする.